\chapter{Assistive Devices}
\index{Assistive Devices}

\section*{Definition and Overview}
Assistive devices are products, equipment, and systems that help people with illness, disability, injury, or ageing to perform everyday activities and to live more independently. They range from simple items such as walking sticks, grab rails, and reading glasses to complex technology such as wheelchairs, hearing aids, communication aids, and prosthetic limbs. The purpose of an assistive device is to compensate for reduced function, to improve safety, and to maintain or restore independence and quality of life. The World Health Organization uses the term assistive technology to describe the broader field of devices and services that support people with functional limitations.

\section*{Epidemiology}
The need for assistive devices is widespread and growing as populations age and as more people live with chronic conditions and disability. The World Health Organization estimates that more than one in ten people worldwide needs at least one assistive product, a figure expected to rise substantially over coming decades. Older adults are the largest group of users, most commonly needing mobility aids, vision aids, and hearing aids. There is also a significant global gap between need and access: many more people could benefit from assistive devices than currently receive them, particularly in low- and middle-income countries.

\section*{Aetiology and Pathophysiology}
Assistive devices do not cause disease; they are prescribed because of underlying health problems that limit function. Common reasons for needing devices include:
\begin{itemize}
    \item \textbf{Mobility problems:} from arthritis, stroke, falls, fractures, amputation, or neurological conditions such as Parkinson's disease
    \item \textbf{Sensory loss:} hearing impairment and vision loss
    \item \textbf{Weakness or fatigue:} from chronic heart, lung, or neuromuscular disease
    \item \textbf{Reduced hand function:} from arthritis, stroke, or injury
    \item \textbf{Communication difficulties:} from stroke, speech disorders, or developmental disability
    \item \textbf{Ageing:} a gradual decline in strength, balance, vision, and hearing
    \item \textbf{Recovery:} temporary devices, such as crutches or walking frames, after surgery or injury
\end{itemize}

\section*{Clinical Features}
The need for an assistive device is usually recognised through difficulties with daily activities:
\begin{itemize}
    \item Trouble walking, climbing stairs, or getting up from a chair
    \item Unsteadiness, or a history of falls and near-falls
    \item Difficulty reaching, gripping, or manipulating objects
    \item Problems hearing conversation or seeing clearly enough for daily tasks
    \item Difficulty with dressing, bathing, toileting, or preparing meals
    \item Reduced endurance that limits outings and social participation
    \item Carer strain from helping with transfers and personal care
\end{itemize}

\section*{Differential Diagnosis}
When assessing the need for a device, the underlying problem must be identified so that the right device is selected.
\begin{itemize}
    \item Balance problems versus muscle weakness versus joint pain as the cause of unsteadiness
    \item A single correctable problem, such as a cataract or ear wax, versus permanent sensory loss
    \item Temporary versus permanent disability
    \item A condition that needs treatment or therapy before a device becomes useful
    \item Whether the problem is best met by a device, a home modification, or a change in activity
\end{itemize}

\section*{When to Seek Medical Care}
See a doctor, physiotherapist, or occupational therapist if mobility, safety, or daily function is declining, if you have had falls, or if you are unsure which device is appropriate. Seek urgent care after any fall that causes significant pain, an inability to bear weight, a blow to the head, or loss of consciousness.

\textbf{Contact your local emergency services for:}
\begin{itemize}
    \item Sudden weakness, numbness, or difficulty speaking, which may indicate a stroke
    \item A fall with severe pain or a suspected fracture
    \item Chest pain, breathlessness, or collapse
    \item Any sudden, severe symptom that causes concern
\end{itemize}

\section*{Diagnosis and Investigations}
\begin{itemize}
    \item \textbf{Functional assessment:} a physiotherapist or occupational therapist assesses strength, balance, mobility, and the way daily activities are managed
    \item \textbf{Medical assessment:} review of underlying conditions and medications that may affect balance, strength, or cognition
    \item \textbf{Vision and hearing tests:} performed by an optometrist or audiologist where sensory loss is a factor
    \item \textbf{Home assessment:} identification of hazards and suitable locations for rails and home modifications
    \item \textbf{Equipment trial:} trying devices to ensure correct fit, height, and ease of use before they are supplied
\end{itemize}

\section*{Management}
\begin{itemize}
    \item \textbf{Mobility aids:} walking sticks, walking frames, crutches, wheelchairs, and mobility scooters
    \item \textbf{Sensory aids:} spectacles, magnifiers, hearing aids, and alerting devices
    \item \textbf{Bathing and toileting aids:} shower chairs, raised toilet seats, and handrails
    \item \textbf{Personal care aids:} dressing and grooming aids for people with reduced hand function
    \item \textbf{Communication and cognitive aids:} speech-generating devices, alarm systems, and memory aids
    \item \textbf{Prostheses and orthoses:} artificial limbs, splints, and supportive footwear
    \item \textbf{Home modifications:} ramps, stair rails, grab bars, and improved lighting
    \item \textbf{Training and review:} instruction in the correct and safe use of the device, with follow-up to adjust it as needs change
\end{itemize}

\section*{Complications}
\begin{itemize}
    \item Falls caused by using an unsuitable, badly fitted, or poorly maintained device
    \item Pressure sores from poorly fitted wheelchairs, cushions, or prostheses
    \item Skin rubbing, blisters, or nerve pressure from splints and supports
    \item Over-dependence on a device, which may reduce strength and physical activity
    \item Tripping hazards created by crutches, frames, or walkers in confined spaces
    \item Devices that are damaged, lost, or left unused because of inadequate maintenance or training
\end{itemize}

\section*{Prognosis and Outcomes}
Assistive devices, when well matched to the person and their environment, improve safety, independence, and quality of life, and may delay or avoid the need for residential care. Outcomes depend on correct assessment, fitting, and training. Some people need a device only during recovery and later stop using it, while others require long-term or progressively increasing support as their condition evolves. Regular review ensures that the device continues to meet the person's needs and is used safely.

\section*{Prevention and Self-Care}
\begin{itemize}
    \item Have the device professionally prescribed, fitted, and adjusted
    \item Keep devices in good repair, and check brakes, grips, and wheels regularly
    \item Reduce fall risks at home with good lighting, clear walkways, and grab rails
    \item Keep active with strengthening and balance exercises to maintain function
    \item Use devices as trained, rather than abandoning or improvising with them
    \item Review devices regularly as needs, weight, or health change
\end{itemize}

\section*{Sources and Further Reading}
The following organisations provide reliable information on assistive devices for patients, carers, and healthcare students.

\begin{itemize}
    \item World Health Organization. \textit{Assistive technology and devices.} https://www.who.int/health-topics/assistive-technology
    \item NHS. \textit{Equipment for daily living and disability support.} https://www.nhs.uk/conditions/social-care-and-support-guide/
    \item MedlinePlus. \textit{Assistive devices and mobility aids.} https://medlineplus.gov/mobilityaids.html
    \item Centers for Disease Control and Prevention (CDC). \textit{Assistive technology and disability.} https://www.cdc.gov/disabilities/
    \item Mayo Clinic. \textit{Independent living and assistive equipment.} https://www.mayoclinic.org/healthy-lifestyle/healthy-aging/
    \item National Institute on Aging (NIA). \textit{Assistive devices for older adults.} https://www.nia.nih.gov/health/
\end{itemize}
